\subsection{函数演算的应用}

命题 9. --- 星代数的每个单射态射都是等距的，且特别地有闭像。

设 A 与 A′ 是星代数，\(\pi\colon A\to A'\) 是从 A 到 A' 的一个对合代数态射。

首先设 A 与 A' 都是幺的且 \(\pi\) 是幺态射。我们有 \(\|\pi\| \leqslant 1\)（命题 2）。设存在 A 中的 x 使 \(\|\pi(x)\| < \|x\|\)。令 \(y = x^{*}x\)；它是 A 的一个 Hermite 元。由于 A 与 A' 都是星代数，我们有 \(\|\pi(y)\| = \|\pi(x)\|^{2} < \|x\|^{2} = \|y\|\)，即 \(\varrho(\pi(y)) < \varrho(y)\) (I, p.~104, 引理 7)。特别地，\(\mathrm{Sp}_{\mathrm{A}'}(\pi(y))\) 是 \(\mathrm{Sp}_{\mathrm{A}}(y)\) 的一个闭子集，且异于 \(\mathrm{Sp}_{\mathrm{A}}(y)\)

(I, p.~3 的注记 6 与 I, p.~24 的定理 1)。于是存在非零函数 \(f \in \mathcal{C}(\mathrm{Sp}_{\mathrm{A}}(y))\) 使 \(f| \mathrm{Sp}_{\mathrm{A}'}(\pi(y)) = 0\) (TG, IX, p.~13, 命题 3)。令 \(w = f(y) \in \mathrm{A}\)。我们有 \(w \neq 0\)，因为 \(f \neq 0\)；但 \(\pi(w) = \pi(f(y)) = f(\pi(y)) = 0\)，因为 \(f\) 在 \(\mathrm{Sp}_{\mathrm{A}'}(\pi(y))\) 上为零（命题 8）。故 \(\pi\) 不是单射。

现在处理一般情形。设 \(\widetilde{\mathrm{A}}\) 与 \(\widetilde{\mathrm{A}}'\) 分别是由 A 与 \(\mathrm{A}'\) 添加单位元而得的星代数 (I, p.~106, 定义 4)。存在唯一的幺对合代数态射 \(\widetilde{\pi}\colon \widetilde{\mathrm{A}}\to \widetilde{\mathrm{A}}'\)，它拓展 \(\pi\)。这个态射是单射，因而由前所证是等距的。于是对每个 \(x\in \mathrm{A}\)，有 \(\| \pi (x)\| = \| \widetilde{\pi} (x)\| = \| x\|\)。

引理 11. --- 设 X 是一个完全正则拓扑空间，即可一致化且分离的拓扑空间 (TG, IX, p.~8, 定义 4)，至少含两个点。存在 \(\mathcal{C}(\mathrm{X})\) 中非零的连续函数 f 与 g 使 \(fg = 0\)。

设 \(x \neq y\) 是 X 的两个不同点。设 U 与 V 分别是 x 与 y 的开邻域，使得 U 与 V 不相交。由于 X 是可一致化的，依 TG, IX, p.~7, 定理 2，存在函数 \(f \in \mathcal{C}(X)\) 使 \(f(x) = 1\) 且 \(f|X - U = 0\)。类似地，存在 \(g \in \mathcal{C}(X)\) 使 \(g(y) = 1\) 且 \(g|X - V = 0\)。于是我们有 fg = 0。

命题 10. --- 设 A 是一个幺星代数。假设对 A 的每对可交换元素 \((x,y)\)，条件 xy = 0 蕴含 x = 0 或 y = 0。那么 A = C · 1。

若 A 不等于 C·1，则存在 A 中的 Hermite 元 x 不属于 C·1 (I, p.~96, 引理 2)。设 B 是 A 的由 x 生成的幺星子代数。它是交换的，且同构于 \(\mathcal{C}(\mathrm{Sp}_{\mathrm{A}}(x))\) (I, p.~110, 注记)。由于 x 不是标量，它在 B 中的谱不缩减为单个元素 (I, p.~110, 例 1)。因此存在 \(\mathrm{Sp}_{\mathrm{A}}(x)\) 上非零的连续函数 f 与 g 使 fg = 0（引理 11）。A 的元素 \(f(x)\) 与 \(g(x)\) 非零、可交换，且在 A 中满足 \(f(x)g(x) = 0\)。

命题 11. --- 设 A 是一个幺星代数，a、x 与 y 是 A 的元素。设 x 与 y 是正规元。若 xa = ay，则对 x 的谱与 y 的谱的并上每个连续函数 f，有 \(f(x)a = af(y)\)。特别地，\(f(aa^{*})a = af(a^{*}a)\) 对每个函数 \(f \in \mathcal{C}(\mathrm{Sp}'(a^{*}a))\) 成立。

设 \(\mathrm{S} = \mathrm{Sp}(x) \cup \mathrm{Sp}(y)\)。I, p.~106 的命题 6 蕴含 \(x^{*}a = ay^{*}\)。从而，我们得到 \(f(x)a = af(y)\)，对每个函数 \(f\)（形如 \(z \mapsto \mathrm{P}(z,\overline{z})\)，其中 \(\mathrm{P} \in \mathbf{C}[\mathrm{X},\mathrm{Y}]\) 是一个多项式）。由于由函数 \(f \in \mathcal{C}(\mathrm{S})\)（满足 \(f(x)a = af(y)\)）组成的集合是 \(\mathcal{C}(\mathrm{S})\) 的一个闭子代数，依 TG, X, p.~40, 推论 1，它与 \(\mathcal{C}(\mathrm{S})\) 重合。

第二个断言是第一个断言应用于 Hermite 元 \(x = aa^{*}\) 与 \(y = a^{*}a\) 的推论，其中用到 \(\mathrm{Sp}'(a^{*}a) = \mathrm{Sp}'(aa^{*})\) 这一事实 (I, p.~5, 命题 1)。

